\documentclass[11pt]{article}
\usepackage[T1]{fontenc}
\usepackage[utf8]{inputenc}
\usepackage[margin=1in]{geometry}
\usepackage{amsmath,amssymb,amsthm,booktabs,longtable,array,microtype}
\usepackage[colorlinks=true,allcolors=blue]{hyperref}
\usepackage{xurl}
% Generated by code/make_tables.py; do not edit.
\newcommand{\Zero}{0}
\newcommand{\One}{1}
\newcommand{\Two}{2}
\newcommand{\Three}{3}
\newcommand{\Four}{4}
\newcommand{\Five}{5}
\newcommand{\Six}{6}
\newcommand{\Seven}{7}
\newcommand{\Eight}{8}
\newcommand{\Nine}{9}
\newcommand{\Ten}{10}
\newcommand{\Eleven}{11}
\newcommand{\Twelve}{12}
\newcommand{\Thirteen}{13}
\newcommand{\Fourteen}{14}
\newcommand{\Fifteen}{15}
\newcommand{\Seventeen}{17}
\newcommand{\Nineteen}{19}
\newcommand{\TwentyOne}{21}
\newcommand{\TwentyThree}{23}
\newcommand{\TwentyFour}{24}
\newcommand{\TwentyNine}{29}
\newcommand{\Sixty}{60}
\newcommand{\EightyFour}{84}
\newcommand{\OneHundredFive}{105}
\newcommand{\EightForty}{840}
\newcommand{\Thousand}{1000}
\newcommand{\AIAstra}{GPT-6 Astra}
\newcommand{\AIOpus}{Claude Opus 5.5 (Anthropic)}
\newcommand{\AIPlatform}{the Genspark platform}
\newcommand{\RepoLink}{\url{https://github.com/jpbald93/multibase-artin-correlations}}
\newcommand{\AuthorName}{Joshua Bald}
\newcommand{\Affiliation}{Independent researcher}
\newcommand{\AuthorOrcid}{0009-0002-1317-6489}
\newcommand{\TextLicense}{CC BY 4.0}
\newcommand{\CodeLicense}{MIT}
\newcommand{\CensusBound}{1\,000\,000\,000}
\newcommand{\LargeBound}{10\,000\,000\,000}
\newcommand{\SmallBound}{1\,000\,000}
\newcommand{\PairCount}{50\,847\,530}
\newcommand{\LargePairCount}{455\,052\,507}
\newcommand{\BaseCount}{64}
\newcommand{\SelectedCount}{12}
\newcommand{\PositiveCount}{16}
\newcommand{\NegativeCount}{48}
\newcommand{\BaseMaximum}{105}
\newcommand{\LowerPrime}{7}
\newcommand{\PositiveBases}{15,26,30,35,38,39,42,51,55,66,70,78,87,95,102,105}
\newcommand{\SelectedBases}{2,3,5,6,7,10,11,13,15,17,21,29}
\newcommand{\BaseFifteenDelta}{+0.012196}
\newcommand{\BaseFifteenLargeDelta}{+0.004646}
\newcommand{\RatioMinimum}{0.969}
\newcommand{\RatioMaximum}{1.040}
\newcommand{\RatioDigits}{3}
\newcommand{\OffDiagonalCount}{132}
\newcommand{\NegativeOffDiagonal}{101}
\newcommand{\OffMean}{-0.003454}
\newcommand{\AsymMean}{0.008958}
\newcommand{\AsymMax}{0.051099}
\newcommand{\AsymPair}{6,10}
\newcommand{\DiagonalHits}{9}
\newcommand{\LawCount}{7\,604\,098}
\newcommand{\LawViolations}{0}
\newcommand{\NaivePairs}{78\,494}
\newcommand{\NaiveEligible}{12\,278}
\newcommand{\NaiveViolations}{0}
\newcommand{\EulerCutoff}{10\,000\,000}
\newcommand{\EulerBound}{0.0000001}
\newcommand{\ArtinLower}{0.373955778}
\newcommand{\ArtinUpper}{0.373955816}
\newcommand{\MaxGap}{282}
\newcommand{\MainRuntime}{41.63}
\newcommand{\LargeRuntime}{38.67}
\newcommand{\NaiveRuntime}{4.65}
\newcommand{\PositiveMinConductor}{60}
\newcommand{\PairStartOld}{5}
\newcommand{\DensityPrecision}{40}
\newcommand{\ValidationBound}{1000}
\newcommand{\ValidationBlocks}{2,17,31,1000}
\newcommand{\Threads}{28}
\newcommand{\ChannelSlots}{150}
\newcommand{\GapCapacity}{300}
\newcommand{\EndpointPair}{5,7}
\newcommand{\FreshExceptionPair}{7,11}
\newcommand{\OldCutoff}{100\,000}

\newtheorem{proposition}{Proposition}
\newcommand{\A}{\mathcal A}
\newcommand{\E}{\mathbb E}
\newcommand{\Cov}{\operatorname{Cov}}
\newcommand{\Var}{\operatorname{Var}}
\title{Consecutive-prime Artin correlations across squarefree bases:\ a census, a residue-channel reconstruction,\ and lagged cross-base correlations}
\author{\AuthorName\\\Affiliation\\\small ORCID: \AuthorOrcid}
\date{}
\begin{document}
\maketitle
\begin{abstract}
This data note records freshly computed primitive-root indicators on consecutive
prime pairs. A census over the \BaseCount\ squarefree bases from $\Two$ through
$\BaseMaximum$, with $\LowerPrime\le p<q<\CensusBound$, gives
\PositiveCount\ positive consecutive-pair associations; the remainder are negative.
For base $\Fifteen$, extending the computation to $\LargeBound$ gives
$\delta=\BaseFifteenLargeDelta$, compared with $\BaseFifteenDelta$ at the smaller bound.
For a selected set of \SelectedCount\ bases, we release the full lagged cross-base
matrix and reconstruct the same-base associations by combining conditional Artin
progression densities with measured residue--gap channel frequencies. The ratio of
reconstructed to empirical association lies between \RatioMinimum\ and
\RatioMaximum. These frequencies are inputs: this is an empirical-weight
consistency check, not an independent prediction. An elementary quadratic-character
argument gives a worked example of a mixed-base obstruction for bases $\Two$ and $\Six$ in either
order at gaps congruent to $\Ten$ or $\Fourteen$ modulo $\TwentyFour$.
The numerical results are computations; the displayed obstruction is an elementary
proof within an existing framework. No asymptotic claim is made.
\end{abstract}

\section{Scope and relationship to earlier work}
For an integer base $a$ and a prime $p$, write $\A_a(p)$ for the indicator that
$a$ is a primitive root modulo $p$; it is zero when $p$ divides $a$.
Artin's conjecture concerns the density of primes for which this indicator is one.
Hooley proved the expected density under the relevant generalized Riemann
hypotheses (GRH)~\cite{hooley}; Moree surveys the density formulas, their
entanglement corrections, and primitive roots in progressions~\cite{moree}.
Here the question is finite and empirical: how do the indicators at adjacent
primes vary across bases, and how much of their association is reproduced by
single-prime progression rates combined with measured prime transitions?

The distinction between a single-prime density and a consecutive-pair law is
essential. Lemke Oliver and Soundararajan give a conjectural quantitative
framework for biases in successive reduced residue classes~\cite{los}.
Their Hardy--Littlewood-based framework does not make the present measured
weights into theoretical inputs. Nor does GRH for the individual Artin densities
supply conditional independence for adjacent prime indicators.

Bald's consolidated manuscript~\cite{baldconsolidated} treats the base-$\Ten$
consecutive-prime results, their exact residue decomposition, same-prime
cross-base correlations, and computational usefulness. We refer there rather
than reproduce those analyses. Bald's quadratic exclusion note~\cite{baldexclusion}
provides the same-base setting. The additions here are a broader squarefree-base
census, a uniform multi-base reconstruction, and a matrix in which the two bases
are attached to \emph{different consecutive primes}. These are dataset extensions,
not a proposed theory of correlations.

Tinkov\'{a}, Waxman and Zindulka study Artin twin primes~\cite{twz}.
Their discussion already allows distinct primitive-root bases at distinct shifted
primes. Thus neither the mixed-base framework nor the necessary quadratic
nonresidue condition is claimed as new. The small worked example below is
included to make the arithmetic behind particular matrix entries transparent.

\section{Definitions and exact finite conventions}
For a cutoff $x$, let $\mathcal P_x$ consist of consecutive rational prime pairs
$(p,q)$ with $\LowerPrime\le p<q<x$. Consecutiveness is in the full rational
prime sequence, not in a subsequence selected by primitive-root status or by a
residue condition. Both endpoints must be below the bound. Give all pairs equal
weight. For a source base $a$ and a destination base $b$, put
\[
 X=\A_a(p),\qquad Y=\A_b(q),\qquad
 n_{ij}=\#\{(p,q)\in\mathcal P_x:X=i,Y=j\}.
\]
Our directional association is
\begin{align}
 \delta(a\mathbin{\to}b;x)
 &=\Pr(Y=\One\mid X=\One)-\Pr(Y=\One\mid X=\Zero)\label{eq:delta}\\
 &=\frac{n_{\One\One}}{n_{\One\Zero}+n_{\One\One}}
   -\frac{n_{\Zero\One}}{n_{\Zero\Zero}+n_{\Zero\One}}
 =\frac{\Cov(X,Y)}{\Var(X)}.\nonumber
\end{align}
The marginals are nondegenerate in every reported table. For $a=b$, abbreviate
this to $\delta_a(x)$. The normalization uses the \emph{source} variance, not a
symmetric Pearson normalization. In particular, the difference between the two
orders in the matrix can involve both covariance and marginal normalization.

Adjacent pairs overlap at a prime. We report exact finite counts and descriptive
summaries only. There are no independent-pair error bars, null-test conclusions,
or claims that a descriptive score certifies a sign. The two cutoffs for base
$\Fifteen$ are nested cumulative samples, not independent replications.

The census bases are precisely the squarefree integers in
$[\Two,\BaseMaximum]$; this is not the set of all nonsquare integers in that
interval. The selected matrix and reconstruction bases are
\[
 \mathcal B=\{\SelectedBases\}.
\]
For squarefree positive $a$, the fundamental discriminant and quadratic
conductor are
\[
 D_a=\begin{cases}a,&a\equiv\One\pmod{\Four},\\
 \Four a,&\text{otherwise},\end{cases}
 \qquad \chi_a(r)=\left(\frac{D_a}{r}\right).
\]
The symbol is the Kronecker symbol, and on odd primes not dividing $a$ it agrees
with the Legendre symbol of $a$. This convention covers the even conductors
without treating the character as a Legendre symbol at a composite denominator.

\section{The squarefree-base census}
There are $\PairCount$ pairs in the main census. Table~\ref{tab:census} gives all
contingency counts, conductors and associations, allowing the displayed decimal
values to be checked without rerunning a prime sieve. Exactly \PositiveCount\
of the \BaseCount\ associations are positive and \NegativeCount\ are negative.
The positive bases are
\[
 \{\PositiveBases\}.
\]
Within this finite set their minimum quadratic conductor is
$\PositiveMinConductor$. This is a description of the recorded sample, not a
sign criterion or a rule for bases outside it. No fitted character statistic or
sign predictor is used.

The positive values show why base-$\Ten$ anticorrelation should not be silently
extended to every base at the same bound. Conversely, positivity at this bound
says nothing about an eventual sign. A conductor is useful metadata, not by
itself an explanation of the ordering of all the measured values.

For base $\Fifteen$, the fresh extension has $\LargePairCount$ pairs and gives
\[
 \delta_{\Fifteen}(\CensusBound)=\BaseFifteenDelta,
 \qquad
 \delta_{\Fifteen}(\LargeBound)=\BaseFifteenLargeDelta.
\]
It remains positive and is smaller at the larger cutoff. These two values do
not determine a decay rate or limiting behavior. In particular, treating the
larger nested sample as a new independent positive-sign experiment would
misrepresent how the data were obtained.

\begingroup\small
\setlength{\tabcolsep}{4pt}
\begin{longtable}{r r r r r r r}
\caption{Squarefree-base census. All consecutive pairs have $\LowerPrime\le p<q<\CensusBound$. Counts use the order $(X,Y)$; $D_a$ is the quadratic conductor.}\label{tab:census}\\
\toprule $a$ & $D_a$ & $n_{\Zero\Zero}$ & $n_{\Zero\One}$ & $n_{\One\Zero}$ & $n_{\One\One}$ & $\delta_a$\\\midrule\endfirsthead
\toprule $a$ & $D_a$ & $n_{\Zero\Zero}$ & $n_{\Zero\One}$ & $n_{\One\Zero}$ & $n_{\One\One}$ & $\delta_a$\\\midrule\endhead
\midrule\multicolumn{7}{r}{Continued on next page}\\\endfoot
\bottomrule\endlastfoot
2 & 8 & 19\,331\,100 & 12\,501\,645 & 12\,501\,645 & 6\,513\,140 & -0.050199\\
3 & 12 & 19\,373\,402 & 12\,459\,637 & 12\,459\,638 & 6\,554\,853 & -0.046676\\
5 & 5 & 17\,887\,263 & 12\,944\,550 & 12\,944\,551 & 7\,071\,166 & -0.066563\\
6 & 24 & 19\,630\,351 & 12\,203\,765 & 12\,203\,765 & 6\,809\,649 & -0.025205\\
7 & 28 & 19\,766\,161 & 12\,065\,063 & 12\,065\,063 & 6\,951\,243 & -0.013491\\
10 & 40 & 19\,758\,045 & 12\,072\,868 & 12\,072\,869 & 6\,943\,748 & -0.014140\\
11 & 44 & 19\,700\,431 & 12\,129\,815 & 12\,129\,814 & 6\,887\,470 & -0.018909\\
13 & 13 & 19\,267\,755 & 12\,442\,700 & 12\,442\,700 & 6\,694\,375 & -0.042573\\
14 & 56 & 19\,839\,940 & 11\,992\,198 & 11\,992\,198 & 7\,023\,194 & -0.007390\\
15 & 60 & 20\,071\,423 & 11\,759\,377 & 11\,759\,377 & 7\,257\,353 & +0.012196\\
17 & 17 & 19\,445\,115 & 12\,315\,810 & 12\,315\,811 & 6\,770\,794 & -0.033025\\
19 & 76 & 19\,879\,965 & 11\,951\,632 & 11\,951\,633 & 7\,064\,300 & -0.003971\\
21 & 21 & 19\,586\,383 & 12\,336\,717 & 12\,336\,717 & 6\,587\,713 & -0.038345\\
22 & 88 & 19\,862\,731 & 11\,969\,812 & 11\,969\,811 & 7\,045\,176 & -0.005518\\
23 & 92 & 19\,891\,024 & 11\,939\,021 & 11\,939\,021 & 7\,078\,464 & -0.002878\\
26 & 104 & 19\,948\,674 & 11\,884\,698 & 11\,884\,699 & 7\,129\,459 & +0.001614\\
29 & 29 & 19\,630\,954 & 12\,177\,294 & 12\,177\,294 & 6\,861\,988 & -0.022422\\
30 & 120 & 20\,244\,448 & 11\,589\,080 & 11\,589\,080 & 7\,424\,922 & +0.026445\\
31 & 124 & 19\,906\,368 & 11\,926\,243 & 11\,926\,243 & 7\,088\,676 & -0.001859\\
33 & 33 & 19\,755\,228 & 12\,111\,764 & 12\,111\,764 & 6\,868\,774 & -0.018187\\
34 & 136 & 19\,897\,813 & 11\,932\,677 & 11\,932\,677 & 7\,084\,363 & -0.002355\\
35 & 140 & 19\,980\,168 & 11\,850\,522 & 11\,850\,522 & 7\,166\,318 & +0.004542\\
37 & 37 & 19\,656\,884 & 12\,159\,789 & 12\,159\,789 & 6\,871\,068 & -0.021134\\
38 & 152 & 19\,934\,198 & 11\,896\,013 & 11\,896\,014 & 7\,121\,305 & +0.000731\\
39 & 156 & 20\,065\,680 & 11\,764\,518 & 11\,764\,518 & 7\,252\,814 & +0.011777\\
41 & 41 & 19\,695\,440 & 12\,122\,551 & 12\,122\,551 & 6\,906\,988 & -0.018035\\
42 & 168 & 20\,038\,405 & 11\,794\,912 & 11\,794\,912 & 7\,219\,301 & +0.009158\\
43 & 172 & 19\,910\,105 & 11\,920\,219 & 11\,920\,219 & 7\,096\,987 & -0.001305\\
46 & 184 & 19\,904\,623 & 11\,925\,964 & 11\,925\,964 & 7\,090\,979 & -0.001793\\
47 & 188 & 19\,917\,453 & 11\,915\,193 & 11\,915\,193 & 7\,099\,691 & -0.000932\\
51 & 204 & 20\,005\,626 & 11\,827\,036 & 11\,827\,036 & 7\,187\,832 & +0.006473\\
53 & 53 & 19\,733\,025 & 12\,092\,997 & 12\,092\,996 & 6\,928\,512 & -0.015726\\
55 & 220 & 20\,045\,243 & 11\,785\,063 & 11\,785\,063 & 7\,232\,161 & +0.010049\\
57 & 57 & 19\,863\,814 & 11\,977\,842 & 11\,977\,842 & 7\,028\,032 & -0.006387\\
58 & 232 & 19\,880\,001 & 11\,951\,635 & 11\,951\,635 & 7\,064\,259 & -0.003972\\
59 & 236 & 19\,870\,776 & 11\,963\,818 & 11\,963\,819 & 7\,049\,117 & -0.005058\\
61 & 61 & 19\,765\,583 & 12\,059\,798 & 12\,059\,799 & 6\,962\,350 & -0.012924\\
62 & 248 & 19\,930\,390 & 11\,904\,047 & 11\,904\,046 & 7\,109\,047 & -0.000033\\
65 & 65 & 19\,899\,839 & 11\,937\,009 & 11\,937\,008 & 7\,073\,674 & -0.002854\\
66 & 264 & 20\,023\,849 & 11\,807\,704 & 11\,807\,704 & 7\,208\,273 & +0.008121\\
67 & 268 & 19\,907\,419 & 11\,925\,775 & 11\,925\,775 & 7\,088\,561 & -0.001832\\
69 & 69 & 19\,855\,663 & 11\,981\,295 & 11\,981\,295 & 7\,029\,277 & -0.006577\\
70 & 280 & 20\,050\,133 & 11\,781\,978 & 11\,781\,978 & 7\,233\,441 & +0.010270\\
71 & 284 & 19\,906\,892 & 11\,927\,214 & 11\,927\,214 & 7\,086\,210 & -0.001973\\
73 & 73 & 19\,800\,404 & 12\,026\,529 & 12\,026\,529 & 6\,994\,068 & -0.010162\\
74 & 296 & 19\,921\,196 & 11\,913\,394 & 11\,913\,394 & 7\,099\,546 & -0.000822\\
77 & 77 & 19\,877\,215 & 11\,955\,403 & 11\,955\,403 & 7\,059\,509 & -0.004309\\
78 & 312 & 19\,940\,372 & 11\,892\,936 & 11\,892\,936 & 7\,121\,286 & +0.000924\\
79 & 316 & 19\,897\,841 & 11\,933\,990 & 11\,933\,990 & 7\,081\,709 & -0.002494\\
82 & 328 & 19\,877\,198 & 11\,956\,924 & 11\,956\,924 & 7\,056\,484 & -0.004469\\
83 & 332 & 19\,891\,316 & 11\,937\,178 & 11\,937\,178 & 7\,081\,858 & -0.002691\\
85 & 85 & 19\,926\,792 & 11\,906\,935 & 11\,906\,935 & 7\,106\,868 & -0.000261\\
86 & 344 & 19\,901\,808 & 11\,931\,119 & 11\,931\,119 & 7\,083\,484 & -0.002276\\
87 & 348 & 19\,969\,919 & 11\,861\,682 & 11\,861\,683 & 7\,154\,246 & +0.003585\\
89 & 89 & 19\,804\,042 & 12\,025\,963 & 12\,025\,963 & 6\,991\,562 & -0.010181\\
91 & 364 & 19\,915\,387 & 11\,916\,900 & 11\,916\,899 & 7\,098\,344 & -0.001068\\
93 & 93 & 19\,811\,004 & 12\,026\,822 & 12\,026\,821 & 6\,982\,883 & -0.010420\\
94 & 376 & 19\,907\,194 & 11\,920\,838 & 11\,920\,838 & 7\,098\,660 & -0.001308\\
95 & 380 & 19\,969\,782 & 11\,863\,309 & 11\,863\,309 & 7\,151\,130 & +0.003417\\
97 & 97 & 19\,827\,544 & 12\,003\,768 & 12\,003\,768 & 7\,012\,450 & -0.008344\\
101 & 101 & 19\,824\,663 & 12\,006\,863 & 12\,006\,864 & 7\,009\,140 & -0.008609\\
102 & 408 & 19\,933\,636 & 11\,899\,037 & 11\,899\,037 & 7\,115\,820 & +0.000425\\
103 & 412 & 19\,911\,006 & 11\,924\,104 & 11\,924\,105 & 7\,088\,315 & -0.001733\\
105 & 105 & 20\,055\,906 & 11\,771\,510 & 11\,771\,510 & 7\,248\,604 & +0.011248\\
\end{longtable}

\endgroup

\section{An empirical-weight residue-channel reconstruction}
\subsection{Inputs and single-prime rates}
For each $a\in\mathcal B$, set $M_a=\operatorname{lcm}(D_a,\EightForty)$.
A channel is the pair $(r,g)$, where $r=p\bmod M_a$ and $g=q-p$ is the
\emph{exact} gap. Its measured weight is
\[
 w_a(r,g)=\frac{\#\{(p,q)\in\mathcal P_x:p\equiv r\pmod{M_a},\ q-p=g\}}
 {|\mathcal P_x|}.
\]
All channel frequencies are computed from the same prime stream used for the
census. They do not use the primitive-root labels, but they do use the measured
successor distribution. They must therefore be counted as empirical inputs.

For the squarefree bases considered here, with the conductor contained in the
modulus, the GRH Artin progression-rate expression can be written as
\begin{equation}
 \alpha_a(r)=\mathbf{\One}_{(r,M_a)=\One}\,
 \mathbf{\One}_{\chi_a(r)=-\One}
 \prod_{\substack{\ell\mid M_a,\ \ell>\Two\\r\equiv\One\ (\bmod\ell)}}
 \left(\One-\frac{\One}{\ell}\right)
 \prod_{\ell\nmid M_a}\left(\One-\frac{\One}{\ell(\ell-\One)}\right).
 \label{eq:alpha}
\end{equation}
Both products run over primes. The first product describes the odd-power
obstructions whose congruence condition is already fixed by $r$; the quadratic
condition is written separately. The remaining Euler factors average the
unfixed obstructions. This is the squarefree-base, conductor-containing case of
the classical Artin progression formulas, not a formula asserted for arbitrary
perfect-power bases or arbitrary moduli~\cite{moree}. Averaging~\eqref{eq:alpha} over the reduced classes modulo $M_a$ recovers
Hooley's density for base $a$, including the entanglement correction that
arises when $a\equiv\One\pmod{\Four}$, because $\mathbb Q(\sqrt a)$ then lies in
$\mathbb Q(\zeta_{M_a})$ and the restriction $\chi_a(r)=-\One$ encodes it.

We evaluate the Euler product through $\EulerCutoff$, accumulating with
$\DensityPrecision$ decimal digits. That arithmetic precision is not a claim
about the infinite product. If $A_B$ is the truncated Artin product, then
\[
 A_B\left(\One-\frac{\One}{B}\right)\le A\le A_B,
 \qquad B=\EulerCutoff,
\]
since the sum of the omitted factors' deficits is at most
$\sum_{n>B}\One/[n(n-\One)]=\One/B$. Thus a conservative bound for the
relative omitted-factor error is $\EulerBound$. The saved output contains both
bounds, and the reconstruction propagates them. The quoted ratio range uses
only \RatioDigits\ decimal places; fine agreement should not be read beyond
the precision justified by the product and by the model itself.

\subsection{The contraction, and what it does not test}
Model I replaces the two endpoint statuses, within a channel, by independent
Bernoulli variables having rates $\alpha_a(r)$ and $\alpha_a(r+g)$. Define
\begin{align*}
 u_a&=\sum_{r,g}w_a(r,g)\alpha_a(r),&
 v_a&=\sum_{r,g}w_a(r,g)\alpha_a(r+g),\\
 t_a&=\sum_{r,g}w_a(r,g)\alpha_a(r)\alpha_a(r+g),&
 \delta_{a,\mathrm{rec}}&=\frac{t_a-u_av_a}{u_a(\One-u_a)}.
\end{align*}
The argument of $\alpha_a$ is reduced modulo $M_a$. No empirical Artin marginal
is substituted into this normalization, and no fitted scale or intercept is used.
The empirical comparison is exactly~\eqref{eq:delta} on the same pair set.

\begin{table}[htbp]\centering
\begin{tabular}{rrrrr}
\toprule $a$ & $M_a$ & $\delta_{\rm emp}$ & $\delta_{\rm rec}$ & Ratio\\\midrule
2 & 840 & -0.050199 & -0.049831 & 0.9927\\
3 & 840 & -0.046676 & -0.047998 & 1.0283\\
5 & 840 & -0.066563 & -0.066471 & 0.9986\\
6 & 840 & -0.025205 & -0.024430 & 0.9692\\
7 & 840 & -0.013491 & -0.013454 & 0.9973\\
10 & 840 & -0.014140 & -0.014703 & 1.0398\\
11 & 9\,240 & -0.018909 & -0.018977 & 1.0036\\
13 & 10\,920 & -0.042573 & -0.042830 & 1.0060\\
15 & 840 & +0.012196 & +0.012159 & 0.9970\\
17 & 14\,280 & -0.033025 & -0.033055 & 1.0009\\
21 & 840 & -0.038345 & -0.038245 & 0.9974\\
29 & 24\,360 & -0.022422 & -0.022473 & 1.0023\\
\bottomrule
\end{tabular}

\caption{Model I with empirical residue--gap weights. Ratio means
$\delta_{a,\mathrm{rec}}/\delta_{a,\mathrm{emp}}$, not predictive accuracy on a
held-out sample.}\label{tab:reconstruction}
\end{table}

Table~\ref{tab:reconstruction} gives a ratio range from \RatioMinimum\ to
\RatioMaximum. The positive base-$\Fifteen$ entry is reconstructed as positive,
while the other displayed same-base associations are negative. This agreement
is useful as a consistency check that the measured prime-transition distribution
and classical single-prime progression rates can account for much of the
finite-range pattern across these bases.

It is not a Hardy--Littlewood derivation: the gap and residue weights were not
predicted. It is not a test establishing conditional independence: that
independence was imposed in forming the product of rates. A small aggregate
error can conceal compensating within-channel errors. Nor does a good aggregate
fit bound information carried by the predecessor after conditioning; that is a
different question treated for base $\Ten$ in~\cite{baldconsolidated}.

Although the raw weights retain exact gaps for transparency, this particular
contraction uses a gap only through the destination residue. Summing weights
with the same $(r,r+g\bmod M_a)$ leaves the reconstruction unchanged. The
exact-gap representation is therefore a data organization, not an additional
independence mechanism.

\subsection{Finite nonreduced-prime exceptions}
Formula~\eqref{eq:alpha} assigns zero to nonreduced classes. This is an
asymptotic progression convention: a prime dividing the modulus can still be
an Artin prime for the base. In the reported pair range, the channel
$(r,g)=(\Seven,\Four)$ for base $\Seventeen$ contains a doubly-Artin pair,
namely $(\FreshExceptionPair)$, even though its reconstructed product rate is
zero. The fresh stream logs those exact endpoints and the channel array
independently retains their contribution.

For clarity about the lower endpoint, if the range were enlarged to start at
$\PairStartOld$, the additional pair $(\EndpointPair)$ would be another such
exception for base $\Seventeen$ and also an exception for base $\Three$.
Direct modular-order checks identify all of these finite events. They are not
violations of a character obstruction: they arise from primes dividing the
chosen modulus. We do not discard them from the empirical counts or assert
that every zero-rate channel is empty.

\section{Lagged cross-base correlations}
The matrix has source $a$ on rows and destination $b$ on columns. Tables
\ref{tab:matrixleft} and~\ref{tab:matrixright} give all entries; the raw data
also retain each contingency table, rather than just the rounded associations.
The diagonal agrees with the selected same-base census values. The
\OffDiagonalCount\ off-diagonal entries contain \NegativeOffDiagonal\ negative
values and have mean $\OffMean$.

\begin{table}[htbp]\centering\small
\begin{tabular}{rrrrrrr}
\toprule $a\backslash b$ & 2 & 3 & 5 & 6 & 7 & 10\\\midrule
2 & -0.050199 & -0.003005 & -0.014488 & +0.009426 & -0.002789 & -0.010957\\
3 & -0.002725 & -0.046676 & -0.014747 & -0.002801 & -0.023345 & -0.003055\\
5 & +0.003313 & +0.003365 & -0.066563 & +0.003188 & +0.003217 & +0.003329\\
6 & -0.035420 & -0.002852 & -0.014710 & -0.025205 & -0.002883 & +0.041485\\
7 & -0.002899 & -0.019289 & -0.014534 & -0.002819 & -0.013491 & -0.002846\\
10 & -0.002569 & -0.002930 & -0.014715 & -0.009614 & -0.002842 & -0.014140\\
11 & -0.002527 & -0.020131 & -0.014342 & -0.002732 & -0.003744 & -0.002789\\
13 & -0.000085 & -0.000135 & -0.010909 & +0.000047 & +0.000054 & -0.000118\\
15 & -0.002858 & +0.022208 & -0.014599 & -0.002823 & +0.001448 & -0.002940\\
17 & -0.004351 & -0.004503 & -0.017752 & -0.004346 & -0.004818 & -0.004413\\
21 & -0.010183 & -0.010350 & -0.033040 & -0.009939 & -0.010200 & -0.009731\\
29 & -0.001001 & -0.001568 & -0.008987 & -0.001271 & -0.001206 & -0.001472\\
\bottomrule
\end{tabular}

\caption{Lagged matrix, first block of destination bases.
Entries are $\delta(a\mathbin{\to}b;\CensusBound)$.}\label{tab:matrixleft}
\end{table}
\begin{table}[htbp]\centering\small
\begin{tabular}{rrrrrrr}
\toprule $a\backslash b$ & 11 & 13 & 15 & 17 & 21 & 29\\\midrule
2 & -0.002754 & -0.006972 & -0.002838 & -0.002966 & +0.004450 & -0.003509\\
3 & -0.003210 & -0.007078 & +0.031438 & -0.002868 & +0.004336 & -0.003712\\
5 & +0.003522 & -0.000430 & +0.003255 & +0.000538 & +0.002513 & +0.006611\\
6 & -0.002943 & -0.007193 & -0.002921 & -0.002932 & +0.004245 & -0.003622\\
7 & -0.002139 & -0.007268 & +0.006698 & -0.002756 & +0.004127 & -0.003728\\
10 & -0.002887 & -0.007276 & -0.002888 & -0.002887 & +0.004185 & -0.003521\\
11 & -0.018909 & -0.007341 & -0.000822 & -0.003006 & +0.004503 & -0.003595\\
13 & +0.000053 & -0.042573 & +0.000117 & +0.000048 & +0.003391 & -0.001002\\
15 & -0.000801 & -0.007280 & +0.012196 & -0.002629 & +0.004258 & -0.003518\\
17 & -0.004325 & -0.008888 & -0.004493 & -0.033025 & +0.001414 & -0.005485\\
21 & -0.009991 & -0.016546 & -0.010018 & -0.008659 & -0.038345 & -0.012848\\
29 & -0.001185 & -0.005830 & -0.001066 & -0.001515 & +0.004144 & -0.022422\\
\bottomrule
\end{tabular}

\caption{Lagged matrix, remaining destination bases. Together with
Table~\ref{tab:matrixleft}, this is the full square matrix.}\label{tab:matrixright}
\end{table}

Across unordered pairs of distinct bases, the mean absolute difference between
opposite directions is $\AsymMean$, with maximum $\AsymMax$ at
$\{\AsymPair\}$. These are descriptive summaries of the directional statistic.
The matrix is not the same-prime cross-base matrix: $\A_a(p)$ and $\A_b(q)$
are measured at different primes, whose residue classes have an oriented
successor distribution. The same-prime entanglement discussion in
\cite{baldconsolidated} cannot simply be relabeled to explain these signs.

In particular, an exclusion of joint successes in some gap classes does not
force the global association to be negative. The global quantity also combines
all other classes and the source/destination marginals. The mixed pair studied
next illustrates a local necessary-condition obstruction, not a global sign
law for the matrix. We make no classification of matrix signs in terms of
multiplicative relations among bases.

\section{A worked mixed quadratic obstruction}
For bases $a,b$, let $L=\operatorname{lcm}(\Two,D_a,D_b)$.
An even shift $g\bmod L$ is called \emph{nonvacuously excluded} when there is at
least one residue $r$ for which both $r$ and $r+g$ are coprime to $L$, but no such
residue with
\[
 \chi_a(r)=\chi_b(r+g)=-\One.
\]
The nonvacuity restriction matters: an odd shift modulo an even modulus cannot
join two reduced residues, and must not be counted as an informative exclusion.
These conditions are necessary quadratic conditions for simultaneous primitive
roots, not sufficient conditions for their existence in all remaining classes.

\begin{proposition}
Let $p,q>\Three$ be primes, $q-p\equiv\Ten$ or $\Fourteen\pmod{\TwentyFour}$.
Then it is impossible for $\Two$ to be a primitive root modulo $p$ and $\Six$
to be a primitive root modulo $q$. The same statement holds with the bases
interchanged. Consecutiveness is not needed.
\end{proposition}
\begin{proof}
For a prime greater than $\Three$, the reduced residues modulo $\TwentyFour$
on which $\chi_{\Two}=-\One$ are
$\Five,\Eleven,\Thirteen,\Nineteen$. Adding either specified gap gives the
finite table below. Whenever the successor residue is reduced, its
$\chi_{\Six}$ value is $+\One$, so $\Six$ cannot be a primitive root there.
If it is not reduced, it cannot be the residue of the stipulated prime $q$.
This proves the first order. For the reversed order, apply the same residue
argument to the reversed endpoint residues: negation of the gap interchanges
$\Ten$ and $\Fourteen$ modulo $\TwentyFour$. The argument depends only on
residues, not on which endpoint is larger.
\end{proof}

\begin{table}[htbp]\centering
\begin{tabular}{rrrl}
\toprule $g$ & $r$ & $r+g\pmod{\TwentyFour}$ & Successor condition\\\midrule
10 & 5 & 15 & not a reduced residue\\
10 & 11 & 21 & not a reduced residue\\
10 & 13 & 23 & $\chi_{\Six}=+\One$\\
10 & 19 & 5 & $\chi_{\Six}=+\One$\\
14 & 5 & 19 & $\chi_{\Six}=+\One$\\
14 & 11 & 1 & $\chi_{\Six}=+\One$\\
14 & 13 & 3 & not a reduced residue\\
14 & 19 & 9 & not a reduced residue\\
\bottomrule
\end{tabular}

\caption{The complete residue check used in the mixed obstruction.}
\end{table}

We enumerate all ordered pairs in $\mathcal B$ directly from the character
values over residues. The only off-diagonal pairs with nonvacuous excluded even
classes are $(\Two,\Six)$ and $(\Six,\Two)$, each with the two classes in the
proposition. Exactly \DiagonalHits\ diagonal pairs have classes; the full
diagonal scan is in Table~\ref{tab:exclusions}. This is an exhaustive statement
about the selected finite set, not a classification for arbitrary bases.

\begin{table}[htbp]\centering
\begin{tabular}{rrl}
\toprule Base & Modulus & Nonvacuous excluded even classes\\\midrule
2 & 8 & 4\\
3 & 12 & 4,6,8\\
5 & 10 & 2,8\\
6 & 24 & 8,12,16\\
7 & 28 & 14\\
10 & 40 & 20\\
11 & 44 & 22\\
13 & 26 & \text{none}\\
15 & 60 & 20,30,40\\
17 & 34 & \text{none}\\
21 & 42 & 14,28\\
29 & 58 & \text{none}\\
\bottomrule
\end{tabular}

\caption{Diagonal part of the finite character scan. An empty set is recorded
as none; it does not assert existence of Artin pairs in every shift.}
\label{tab:exclusions}
\end{table}

In the main prime census, $\LawCount$ pairs have one of the two specified gap
classes, and there are $\LawViolations$ simultaneous successes in either
order. Separately, a deliberately naive primitive-root implementation enumerates
multiplicative orders for both bases at every prime in the range through
$\SmallBound$. Among $\NaivePairs$ consecutive pairs it finds $\NaiveEligible$
with the specified classes and $\NaiveViolations$ violations in either order.
This check validates the implementation against the elementary argument; the
finite computation is not the proof of the all-primes statement.

The proposition is also formally verified in the Lean proof assistant with Mathlib. The package
directory \path{lean/} states both orders as
\path{MixedExclusion.not_primitiveRoot_two_six_of_gap} and
\path{MixedExclusion.not_primitiveRoot_six_two_of_gap}, with the gap condition
written as a congruence of residues modulo $\TwentyFour$ and primitive roots
expressed through Mathlib's \path{IsPrimitiveRoot}. The proof uses Euler's
criterion, the second supplementary law, quadratic reciprocity for $\Three$ and
an exhaustive residue check. The checked statements depend only on Lean's
standard axioms, and the build contains no unproved steps. The formalization
covers this proposition only, not the computations.

\section{Computation, validation and reproducibility}
All large enumerations and the residue/reconstruction calculations were run
fresh on the author's multi-core workstation. The sieve source was transferred as exact
base-encoded bytes and its cryptographic checksum checked before compilation. The package
records source checksums, commands, output checksums and runtimes. The main
runs used \Threads\ OpenMP threads, with approximate census runtimes
$\MainRuntime$ seconds for the multi-base bound and $\LargeRuntime$ seconds
for the single-base extension. These are machine-specific execution records,
not algorithmic speed comparisons.

The C census marks composites in independent half-open integer blocks. Prime
labels are computed within blocks and the resulting streams are assembled in
block order, retaining the last prime and its labels across each boundary.
Thus no halo length or assumed maximum prime gap is needed to recover
consecutiveness. Only the initial integers below the prime range are specially
excluded; the first integers of later blocks are sieved normally.

For each prime, distinct factors of $p-\One$ are obtained once. For every base,
the primitive-root test checks that the base is a unit and that
$a^{(p-\One)/\ell}\not\equiv\One\pmod p$ for every prime divisor $\ell$ of
$p-\One$. Modular multiplication uses unsigned double-width integer arithmetic.
The single-base extension uses the same prime-stream construction and root test.
The gap arrays have \ChannelSlots\ slots for even gaps up to $\GapCapacity$;
there is an assertion, not silent overflow clamping. The observed maximum gap
in the main run is $\MaxGap$.

Before scaling, an independent Python implementation trial-divided
integers for primality and enumerated complete multiplicative orbits through
$\ValidationBound$. For block sizes $\ValidationBlocks$, it checked every
prime label, every same-base and mixed-base count, and every nonzero channel
cell. This includes boundaries deliberately placed inside the small prime
range. The same tiny computation on the workstation agreed exactly.
A separate C orbit-enumeration program, checked on the tiny range, supplies
the naive mixed-law check; it does not reuse the factor-of-$p-\One$ criterion.

For the character calculation, the Kronecker routine was checked against
Euler's criterion on small odd primes. The exclusion scan represents reduced
residue sets and negative-character sets as bitsets and intersects cyclic
shifts. It enumerates every even shift and tests nonvacuity before testing
exclusion. The reconstruction reads the saved channel arrays and applies the
printed formulas without an empirical primitive-root-rate substitution.

The numerical TeX macros, tables and bibliographic identifiers are generated
by the script \path{code/make_tables.py}, using saved outputs and explicit metadata.
The raw binary channel files include their base, modulus, gap-slot count and
pair-count headers; their format is documented in the README. All final
numerical displays can therefore be rebuilt without a fresh sieve, while the
source needed for a full resieve is also supplied. Exact counts are primary;
rounding is confined to generated presentation material.

\section{Interpretation and limitations}
The note adds a portable finite census and a multi-base comparison, not a new
mechanism beyond the residue-channel perspective. The same-base positive
examples prevent an overgeneralization of the base-$\Ten$ behavior. The
lagged matrix records an observable different from same-prime entanglement.
The reconstruction separates measured prime-transition frequencies from
classical single-prime rates, making its empirical inputs explicit.

None of these observations settles the large-cutoff behavior of actual Artin
indicators. In the separable fixed-modulus model considered in
\cite{baldconsolidated}, product limiting residue marginals force the modelled
covariance to vanish: the joint rate is the average of a product of separate
endpoint functions. The present contraction has exactly that separable form.
Unequal finite gap weights do not invalidate that elementary observation.
It is a statement about the restricted model under a limiting residue-law
assumption, not a proof of an asymptotic for the actual indicators. The
finite-range data here do not settle that asymptotic question.

A genuine prediction would require weights obtained without the measured
successor stream, and a justified account of any residual endpoint dependence.
Increasing a table's size or matching its aggregate covariance does not provide
either ingredient. We regard the present package as a reproducible data and
consistency-check supplement to the cited work, with the mixed obstruction as
an explanatory example.

\section*{Data availability and licences}
The reproduction package \texttt{reproduction.zip} accompanies this preprint.
It contains the raw outputs, source programs, generated tables, manuscript,
build instructions and cryptographic checksum manifests. The same material is
public at \RepoLink. No deposit identifier is asserted here.
Text is licensed under \TextLicense; code under \CodeLicense.

\section*{AI assistance}
Large-language-model assistants, accessed through \AIPlatform, were used as
follows. \AIAstra\ wrote the computation and analysis code, ran the
computations, searched the literature and drafted the manuscript. \AIOpus\
coordinated the work, reviewed the draft and carried out a separate
adversarial audit of the mathematics, the numbers and the reproduction
package. That audit included an independent re-implementation of the sieve and
primitive-root test, whose counts matched those reported here. Claude Opus also
wrote the Lean formalization of the mixed obstruction. Every number in the
paper was produced by executing the supplied programs; none was generated by a
language model. The author directed the work and is responsible for the
content.

\begin{thebibliography}{99}
\bibitem{hooley} C. Hooley. On Artin's conjecture. Journal f\"ur die reine und angewandte Mathematik 225 (1967), 209--220. \url{https://doi.org/10.1515/crll.1967.225.209}.
\bibitem{moree} P. Moree. Artin's primitive root conjecture---a survey. Integers 12 (2012), no. 6, 1305--1416. \url{https://doi.org/10.1515/integers-2012-0043}.
\bibitem{los} R. J. Lemke Oliver and K. Soundararajan. Unexpected biases in the distribution of consecutive primes. Proceedings of the National Academy of Sciences 113 (2016), no. 31, E4446--E4454. \url{https://doi.org/10.1073/pnas.1605366113}.
\bibitem{twz} M. Tinkov\'{a}, E. Waxman and M. Zindulka. Artin twin primes. Journal of Number Theory 245 (2023), 203--232. \url{https://doi.org/10.1016/j.jnt.2022.10.006}.
\bibitem{baldexclusion} J. Bald. Quadratic exclusion laws for consecutive Artin primes in arbitrary bases. Zenodo preprint (2026). \url{https://doi.org/10.5281/zenodo.22865343}.
\bibitem{baldconsolidated} J. Bald. Correlations of Artin status among primes: consecutive-prime anticorrelation, cross-base entanglement, and an elementary decomposition. Repository manuscript. \url{https://github.com/jpbald93/artin-correlations}.
\end{thebibliography}

\end{document}
